\begin{tikzpicture}
	
	\colorlet{genotype}{gray}
	\colorlet{phenotype}{OkabeIto4}
	

	%%% SlCLV3 promoter bashing %%%
	\coordinate (SlCLV3 bashing) at (0, 0);
	
	\def\WT{WT}
	\expandafter\def\csname CLV3p13\endcsname{\textit{slclv3\textsuperscript{pro-13}}}
	\expandafter\def\csname CLV3p15\endcsname{\textit{slclv3\textsuperscript{pro-15}}}
	\expandafter\def\csname CLV3p24\endcsname{\textit{slclv3\textsuperscript{pro-24}}}
	\expandafter\def\csname CLV3p26\endcsname{\textit{slclv3\textsuperscript{pro-26}}}
	
	\pgfplotsset{
		gRNA/.style = {arrows = -{Triangle[length = .9pt, width = .2cm]}, OkabeIto2},
		DNA/.style = {line width = .1cm, lightgray, -},
		deletion/.style = {line width = .05cm, OkabeIto7, densely dashed, -},
		xaxis/.is choice,
		xaxis/genotype/.style = {
			xtick = {-2.1, -1.5, -1, -.5},
			extra x ticks = {0},
			extra x tick labels = {ATG},
			xticklabel style = {name = xticklabel},
			xmin = -2.1,
			xmax = 0,
			x decimals = 1,
		},
		xaxis/phenotype/.style = {
			xmin = -3.5
		}		
	}
	
	\begin{pggfplot}[%
		at = {(SlCLV3 bashing)},
		total width = \twocolumnwidth,
		height = 2cm,
		xlabel = {\vphantom{position relative to}\\\vphantom{\textit{ZmCLE7} ATG (kb)}},
		ylabel = {\textit{SlCLV3}\\promoter allele},
		ylabel style = {xshift = -2cm/12},
		ylabel add to width = 4\baselineskip,
		label each x axis,
		enlargelimits = 0,
		ytick = {1, ..., 5},
		yticklabels are csnames,
		title color from name,
		axis background/.style = {fill = none},
		empty facet warning = false,
		tight y
	]{SlCLV3_alleles}
	
		\pggf_quiver(
			data = SlCLV3_alleles_DNA,
			DNA
		)
		
		\pggf_quiver(
			data = SlCLV3_alleles_deletion,
			deletion
		)
		
		\pggf_quiver(
			data = SlCLV3_alleles_gRNA,
			gRNA
		)
		
		\pggf_annotate(
			annotation = {
				\coordinate (sep) at (-1.0805, \pggfymax);
			},
			only facet = 1
		)
		
		\pggf_annotate(
			zeroline*,
			only facet = 2
		)
		
		\pggf_boxplot(
			data = SlCLV3_locules,
			fill = OkabeIto7,
			only facet = 2
		)
	
	\end{pggfplot}
	
	\planticon[height = .5cm]{tomatoFruit} at ($(pggfplot c2r1.north east) + (-.3, -.3)$);
	
	% xlabels
	\node[anchor = north, align = center] at (pggfplot c1r1 |- xlabel.north) {position relative to\\\textit{SlCLV3} ATG (kb)};
	\node[anchor = north, align = center] at (pggfplot c2r1 |- xlabel.north) {locule number per fruit};
	
	% save coordinates
	\coordinate (start) at (pggfplot c1r1.west |- xticklabel.south);
	\coordinate (end) at (pggfplot c1r1.east |- xticklabel.south);
	
	
	%%% SlCLV3 Plant STARR-seq %%%
	\coordinate[yshift = -\columnsep] (SlCLV3 STARR) at (SlCLV3 bashing |- xlabel.south);
	
	\tikzset{
		ROIstart/.store in = \ROIstart,
		ROIend/.store in = \ROIend
	}

	\begin{pggfplot}[%
		at = {(SlCLV3 STARR)},
		total width = \twocolumnwidth,
		xlabel = {position relative to \textit{SlCLV3} ATG (kb)},
		ylabel = \enrichment,
		xmin = -2.1,
		xmax = 0,
		enlarge x limits = 0,
		xtick = {-2.1, -1.8, ..., -.1},
		extra x ticks = {0},
		extra x tick labels = {ATG},
		x decimals = 1,
		row title = {Plant STARR-seq (\usename{dark})},
		row title color = dark,
		axis background/.style = {fill = none},
		legend style = {at = {(1, 1)}},
		every legend image post/.append style = {very thick, arrows = -{Latex[scale = .67]}},
		legend image width = transformdirectionx(0.17),
		tight y
	]{SlCLV3_PlantSTARRseq}
		
		\pggf_annotate(
			zeroline,
			annotation = {
				\coordinate (sep 2) at (-1.0805, \pggfymin);
				\fill[dark, fill opacity = .1] (\ROIstart, \pggfymin) coordinate (ROI Sl1) rectangle (\ROIend, \pggfymax) coordinate (ROI Sl2);
			},
			manual legend = {
				experiment = {black},
				prediction = {35Senhancer}
			}
		)
	
		\pggf_quiver(
			data = SlCLV3_PlantSTARRseq_prediction,
			35Senhancer,
			very thick,
			arrows = -{Latex[scale = .67]},
		)
	
		\pggf_quiver(
			data = SlCLV3_PlantSTARRseq_experiment,
			black,
			very thick,
			arrows = -{Latex[scale = .67]},
		)
	
	\end{pggfplot}
	
	\planticon[height = .5cm]{dark} at ($(pggfplot c1r1.north west) + (.3, -.3)$);
	
	% connection to previous plot
	\begin{pgfonlayer}{background}
		\draw[gray] (start) to[out = -90, in = 90] (pggfplot c1r1.north west) (end) to[out = -90, in = 90, looseness = .75] (pggfplot c1r1.north east);
		\draw[gray, dashed] (sep) -- (sep |- start) to[out = -90, in = 90] (sep 2 |- pggfplot c1r1.north) -- (sep 2);
	\end{pgfonlayer}
	
	
	
	%%% ZmCLE7 promoter bashing %%%
	\coordinate (ZmCLE7 bashing) at (SlCLV3 bashing -| \textwidth - \twocolumnwidth, 0);
	
	\expandafter\def\csname CLE7p1\endcsname{\textit{CLE7\textsuperscript{CR-pro1}}}
	\expandafter\def\csname CLE7p4\endcsname{\textit{CLE7\textsuperscript{CR-pro4}}}
	\expandafter\def\csname CLE7p6\endcsname{\textit{CLE7\textsuperscript{CR-pro6}}}
	
	\pgfplotsset{
		xaxis/genotype/.style = {
			xtick = {-2, -1.5, ..., -.1},
			extra x ticks = {0},
			extra x tick labels = {ATG},
			xticklabel style = {name = xticklabel},
			xmin = -2.313,
			xmax = 0,
			x decimals = 1,
		},
		xaxis/phenotype/.style = {
			xtick = {0, .2, ..., 1.5},
			x decimals = 1,
			xmin = -.15,
			xmax = 1.39
		}		
	}
	
	\begin{pggfplot}[%
		at = {(ZmCLE7 bashing)},
		total width = \twocolumnwidth,
		height = 2cm,
		xlabel = {\vphantom{position relative to}\\\vphantom{\textit{ZmCLE7} ATG (kb)}},
		ylabel = {\textit{ZmCLE7}\\promoter allele},
		ylabel style = {xshift = -.2cm},
		ylabel add to width = 4\baselineskip,
		label each x axis,
		enlargelimits = 0,
		ytick = {1, ..., 4},
		yticklabels are csnames,
		title color from name,
		axis background/.style = {fill = none},
		empty facet warning = false,
		tight y
	]{ZmCLE7_alleles}
	
		\pggf_quiver(
			data = ZmCLE7_alleles_DNA,
			DNA
		)
		
		\pggf_quiver(
			data = ZmCLE7_alleles_deletion,
			deletion
		)
		
		\pggf_quiver(
			data = ZmCLE7_alleles_gRNA,
			gRNA
		)
		
		\pggf_annotate(
			annotation = {
				\coordinate (sep) at (-0.6665, \pggfymax);
				\coordinate (sep 2) at (-1.5005, \pggfymax);
			},
			only facet = 1
		)
		
		\pggf_bar(
			data = ZmCLE7_expression,
			xbar,
			fill = OkabeIto5,
			draw = black,
			x error = {se},
			error bars/error bar style = {black, mark size = 2.5},
			only facet = 2
		)
		
		\pggf_annotate(
			zeroline*,
			only facet = 2
		)
		
		\pggf_scatter(
			data = ZmCLE7_expression,
			y filter/.expression = {(y == 3 && x == 0.498057387) ? y - .15 : ((y == 3 && x == 0.499394362) ? y + .15 : y)},
			only facet = 2
		)
	
	\end{pggfplot}
	
	\planticon[height = .5cm]{maize} at ($(pggfplot c2r1.north east) + (-.3, -.3)$);
	
	% xlabels
	\node[anchor = north, align = center] at (pggfplot c1r1 |- xlabel.north) {position relative to\\\textit{ZmCLE7} ATG (kb)};
	\node[anchor = north, align = center] at (pggfplot c2r1 |- xlabel.north) {\textit{ZmCLE7} expression};
	
	% save coordinates
	\coordinate (start) at (pggfplot c1r1.west |- xticklabel.south);
	\coordinate (start top) at (pggfplot c1r1.west);
	\coordinate (end) at (pggfplot c1r1.east |- xticklabel.south);
	
	
	%%% ZmCLE7 %%%
	\coordinate (ZmCLE7 plantGREP) at (SlCLV3 STARR -| ZmCLE7 bashing);
	
	\begin{pggfplot}[%
		at = {(ZmCLE7 plantGREP)},
		total width = \twocolumnwidth,
		xlabel = {position relative to \textit{ZmCLE7} ATG (kb)},
		ylabel = \enrichment,
		xmin = -2.313,
		xmax = 0,
		enlarge x limits = 0,
		xtick = {-2.4, -2.1, ..., -.1},
		extra x ticks = {0},
		extra x tick labels = {ATG},
		x decimals = 1,
		row title = {plantGREP (\usename{maize})},
		row title color = maize,
		axis background/.style = {fill = none},
		tight y
	]{ZmCLE7_prediction_points}
	
		\pggf_annotate(
			zeroline,
			annotation = {
				\coordinate (sep 3) at (-0.6665, \pggfymin);
				\coordinate (sep 4) at (-1.5005, \pggfymin);
				\fill[maize, fill opacity = .15] (\ROIstart, \pggfymin) coordinate (ROI Zm1) rectangle (\ROIend, \pggfymax) coordinate (ROI Zm2);
			}
		)
		
		\pggf_line(
			data = ZmCLE7_prediction_lines,
			35Senhancer,
			thick,
			on layer = axis foreground
		)
		
		\pggf_scatter(
			35Senhancer,
			opacity = .15
		)
	
	\end{pggfplot}
	
	\planticon[height = .5cm]{maize} at ($(pggfplot c1r1.north west) + (.3, -.3)$);
	
	% connection to previous plot
	\begin{pgfonlayer}{background}
		\draw[gray] (start top) -- (start) to[out = -90, in = 90] (pggfplot c1r1.north west) (end) to[out = -90, in = 90, looseness = .75] (pggfplot c1r1.north east);
		\draw[gray, dashed] (sep) -- (sep |- start) to[out = -90, in = 90, looseness = .75] (sep 3 |- pggfplot c1r1.north) -- (sep 3) (sep 2) -- (sep 2 |- start) to[out = -90, in = 90] (sep 4 |- pggfplot c1r1.north) -- (sep 4);
	\end{pgfonlayer}
	
	
	%%% DeepLIFT predictions for best SlCLV3 fragment %%%
	\coordinate[yshift = -\columnsep] (SlCLV3 DeepLIFT) at (SlCLV3 bashing |- xlabel.south);
	
	\begin{pggfplot}[%
		at = {(SlCLV3 DeepLIFT)},
		total width = \textwidth,
		height = 1.5cm,
		xlabel = {position relative to \textit{SlCLV3} ATG (bp)},
		ylabel = contribution\\score,
		ylabel add to width = 8pt,
		col title = DeepLIFT (\usename{dark}),
		col title color = dark,
		enlarge x limits = 0,
		ytick = {-0.2, -0.1, ..., 0.3},
		xtick = {-1550, -1540, ..., -1350}
	]{SlCLV3_DeepLIFT}
		
		\pggf_annotate(
			zeroline,
			annotation = {
				\fill[dark, fill opacity = .1] (-1503.5, \pggfymin) rectangle (-1498.5, \pggfymax) node[pos = .5, anchor = south, yshift = -.75cm, text = black, opacity = 1, text depth = 0pt, node font = \figsmall] {MYB};
				\fill[dark, fill opacity = .1] (-1492.5, \pggfymin) rectangle (-1480.5, \pggfymax) node[pos = .5, anchor = south, yshift = -.75cm, text = black, opacity = 1, text depth = 0pt, node font = \figsmall] {split bZIP?};
				\fill[dark, fill opacity = .1] (-1471.5, \pggfymin) rectangle (-1465.5, \pggfymax) node[pos = .5, anchor = south, yshift = -.75cm, text = black, opacity = 1, text depth = 0pt, node font = \figsmall] {MYB?};
				\fill[dark, fill opacity = .1] (-1432.5, \pggfymin) rectangle (-1427.5, \pggfymax) node[pos = .5, anchor = south, yshift = -.75cm, text = black, opacity = 1, text depth = 0pt, node font = \figsmall] {MYB};
				\fill[dark, fill opacity = .1] (-1412.5, \pggfymin) rectangle (-1405.5, \pggfymax) node[pos = .5, anchor = south, yshift = -.75cm, text = black, opacity = 1, text depth = 0pt, node font = \figsmall] {bZIP};
				\fill[dark, fill opacity = .1] (-1392.5, \pggfymin) rectangle (-1387.5, \pggfymax) node[pos = .5, anchor = south, yshift = -.75cm, text = black, opacity = 1, text depth = 0pt, node font = \figsmall] {MYB};
				\fill[dark, fill opacity = .1] (-1376.5, \pggfymin) rectangle (-1370.5, \pggfymax) node[pos = .5, anchor = south, yshift = -.75cm, text = black, opacity = 1, text depth = 0pt, node font = \figsmall] {MYB};
			}
		)
	
		\pggf_logo(letter colors = DNA)
	
	\end{pggfplot}
	
	\planticon[height = .5cm]{dark} at ($(pggfplot c1r1.north west) + (.3, -.3)$);
	
	\begin{pgfonlayer}{background}
		\draw[gray] (ROI Sl1) -- (pggfplot c1r1.above north west) (ROI Sl2 |- ROI Sl1) -- (pggfplot c1r1.above north east);
	\end{pgfonlayer}
	

	%%% DeepLIFT predictions for best ZmCLE7 region %%%
	\coordinate[yshift = -\columnsep] (ZmCLE7 DeepLIFT) at (SlCLV3 bashing |- xlabel.south);
	
	\begin{pggfplot}[%
		at = {(ZmCLE7 DeepLIFT)},
		total width = \textwidth,
		height = 1.5cm,
		xlabel = {position relative to \textit{ZmCLE7} ATG (bp)},
		ylabel = contribution\\score,
		ylabel add to width = 8pt,
		col title =  DeepLIFT (\usename{maize}),
		col title color = maize,
		enlarge x limits = 0,
		ytick = {-0.1, -0.05, ..., 0.15},
		xtick = {-550, -540, ..., -350}
	]{ZmCLE7_DeepLIFT}
	
		\pggf_annotate(
			zeroline,
			annotation = {
				\fill[maize, fill opacity = .15] (-487.5, \pggfymin) rectangle (-480.5, \pggfymax) node[pos = .5, anchor = south, yshift = -.75cm, text = black, opacity = 1, text depth = 0pt, node font = \figsmall] {bZIP};
				\fill[maize, fill opacity = .15] (-428.5, \pggfymin) rectangle (-424.5, \pggfymax) node[pos = .5, anchor = south, yshift = -.75cm, text = black, opacity = 1, text depth = 0pt, node font = \figsmall] {GATA};
				\fill[maize, fill opacity = .15] (-409.5, \pggfymin) rectangle (-405.5, \pggfymax) node[pos = .5, anchor = south, yshift = -.75cm, text = black, opacity = 1, text depth = 0pt, node font = \figsmall] {GATA};
				\fill[maize, fill opacity = .15] (-403.5, \pggfymin) rectangle (-394.5, \pggfymax) node[pos = .5, anchor = north, yshift = .75cm, text = black, opacity = 1, text depth = 0pt, node font = \figsmall] {2 $\times$ C2H2};
				\fill[maize, fill opacity = .15] (-364.5, \pggfymin) rectangle (-360.5, \pggfymax) node[pos = .5, anchor = south, yshift = -.75cm, text = black, opacity = 1, text depth = 0pt, node font = \figsmall] {GATA};
			}
		)
	
		\pggf_logo(letter colors = DNA)
	
	\end{pggfplot}
	
	\planticon[height = .5cm]{maize} at ($(pggfplot c1r1.north west) + (.3, -.3)$);
	
	\begin{pgfonlayer}{background}
		\draw[maize!50] (ROI Zm1) -- (pggfplot c1r1.above north west) (ROI Zm2 |- ROI Zm1) -- (pggfplot c1r1.above north east);
	\end{pgfonlayer}
	
	
	%%% subfigure labels %%%
	\subfiglabel{SlCLV3 bashing}
	\subfiglabel{SlCLV3 STARR}
	\subfiglabel{ZmCLE7 bashing}
	\subfiglabel{ZmCLE7 plantGREP}
	\subfiglabel{SlCLV3 DeepLIFT}
	\subfiglabel{ZmCLE7 DeepLIFT}

\end{tikzpicture}